\documentclass[11pt]{article}
\usepackage[margin=1.1in]{geometry}
\usepackage{amsmath,amssymb,amsthm,mathtools}
\usepackage{booktabs}
\usepackage{hyperref}
\usepackage{microtype}

\newtheorem{theorem}{Theorem}
\newtheorem{proposition}{Proposition}
\newtheorem{corollary}{Corollary}
\newtheorem{definition}{Definition}
\newtheorem{conjecture}{Conjecture}
\newtheorem{remark}{Remark}

\newcommand{\C}{\mathbb{C}}
\newcommand{\wedgeN}[2]{\wedge^{#1}\mathcal{H}_{#2}}
\newcommand{\ket}[1]{|#1\rangle}

\title{Algorithmic construction of exact extremal states for fermionic moment\\ polytopes: completing $\wedge^4\mathcal{H}_9$}

\author{James Orlando\thanks{Independent researcher, Derry, NH, USA.
\texttt{jamie@orlandonh.com}. ORCID: 0009-0008-3158-771X.}}

\date{July 2026 --- DRAFT v0.7 \\ \small Preprint: \href{https://doi.org/10.5281/zenodo.21313834}{doi:10.5281/zenodo.21313834}}

\begin{document}
\maketitle

\begin{abstract}
We present an algorithmic pipeline that constructs exact extremal states for
fermionic natural-occupation-number (moment) polytopes, and use it to complete
the $\wedge^4\mathcal{H}_9$ polytope of Altunbulak and Klyachko [Comm. Math.
Phys. 282, 287 (2008)]: the two vertices verified in that work ``only
numerically''---$(16,16,16,6,6,6,6,6,6)/21$ and $(20,14,14,14,14,4,4,4,4)/23$---%
now have explicit extremal states with elementary, hand-checkable proofs, so
the completeness argument for the published $\wedge^4\mathcal{H}_{10}$
constraint list is fully analytic. The pipeline decides, vertex by vertex,
whether a one-hop-independent state (equivalently, a root-distinct support in
the sense of Wildberger--Sjamaar and Maciazek--Tsanov) exists---by integer and
mixed-integer programming with solver-independent certificates---and, when it
does not, repairs minimal exchange-coupled ansatz classes by an exact
phase-feasibility criterion, followed by verification in exact arithmetic. The
resulting census shows all vertices of $\wedge^3\mathcal{H}_6$ and
$\wedge^3\mathcal{H}_7$, and $27$ of $38$ vertices of $\wedge^3\mathcal{H}_8$
($22$ integer at the natural denominator), are attainable by such states, while
the second open $\wedge^4\mathcal{H}_9$ vertex is provably not: certified
infeasibility, reproduced on two independent solvers, shows its extremal
states cannot be one-hop independent in any orbital basis. The explicit state we
construct for it carries a single forced phase, $\cos\gamma = 3/(4\sqrt{14})$;
exhaustive enumeration excludes real amplitudes within its ansatz class, and a
numerical orbit search indicates the state admits no real form under any
orbital rotation (maximal antiunitary-symmetry overlap $\approx 0.9332$),
while the existence of some other real extremal state for this vertex remains
open. We describe how the same pipeline is being scaled toward the
one-particle dimensions ($d\sim20$--$40$) required for quantum-chemical
applications.
\end{abstract}

\section{Introduction}

The Pauli exclusion principle bounds fermionic natural occupation numbers
(NONs) by $0 \le \lambda_i \le 1$. Klyachko's solution of the one-body pure
$N$-representability problem \cite{Klyachko2006,AK2008} showed that the
antisymmetry of $N$-fermion wave functions imposes far stronger conditions:
the achievable ordered spectra
$\lambda_1 \ge \cdots \ge \lambda_d$ of the one-body reduced density matrix
(1-RDM) of a pure state in $\wedgeN{N}{d}$ form a convex polytope
$\Pi_{N,d}$, cut out by finitely many \emph{generalized Pauli constraints}
(GPCs). These constraints have measurable consequences: (quasi-)pinning of
occupation spectra to facets of $\Pi_{N,d}$ constrains the structure of the
wave function itself \cite{SGC2013,Liebert2025}, and the polytopes govern
applications from reduced-density-matrix functional theory to quantum
information \cite{Castillo2021,Liebert2025}.

Altunbulak and Klyachko computed $\Pi_{N,d}$ completely for all systems with
$d \le 8$ (and for $\wedge^3\mathcal{H}_9$, $\wedge^3\mathcal{H}_{10}$,
$\wedge^4\mathcal{H}_9$, $\wedge^4\mathcal{H}_{10}$,
$\wedge^5\mathcal{H}_{10}$), supplying for each vertex of the small systems
an explicit extremal state ``for those who don't trust computer assisted
proofs'' \cite{AK2008}. For $\wedge^4\mathcal{H}_9$, however, two vertices
resisted their representation-theoretic constructions:
\begin{equation}
\label{eq:AB}
v_A = \tfrac{1}{21}(16,16,16,6,6,6,6,6,6),
\qquad
v_B = \tfrac{1}{23}(20,14,14,14,14,4,4,4,4),
\end{equation}
which were, in their words, ``checked only numerically''---and on which the
completeness of the published $\wedge^4\mathcal{H}_{10}$ constraint list
depends \cite{AK2008}. To our knowledge these verifications have remained
numerical for the intervening eighteen years.

\paragraph{Contributions.}
\begin{enumerate}
\item \textbf{An explicit extremal state for $v_A$}
(Theorem~\ref{thm:psiA}), with an elementary proof: seven Slater
determinants, pairwise sharing at most two orbitals, with squared amplitudes
$k_T/21$ for integer weights $k_T$ summing to $21$. The state would occupy a
single line of Table~6 of Ref.~\cite{AK2008}.
\item \textbf{A design/interference dichotomy.} Abstracting the structure of
$\psi_A$, we define \emph{weighted-design} states
(Definition~\ref{def:design}) and show that design-attainability of a vertex
is a pure integer-programming question. An exhaustive census
(Section~\ref{sec:census}) over all vertices of $\Pi_{3,6}$, $\Pi_{3,7}$,
$\Pi_{3,8}$ finds designs for $4/4$, $10/10$, and $22/38$ vertices
respectively; for $v_B$ we prove by mixed-integer infeasibility that
\emph{no} design exists, so every extremal state of $v_B$ requires phase
cancellation between one-hop-connected configurations. The vertices that
resisted the constructions of Ref.~\cite{AK2008} are precisely of this
interference type.
\item \textbf{Canonical forms of GPC polytopes.} We compute
$f$-vector data and the degrees of the adjoint hypersurfaces (equivalently,
of the numerators of the Arkani-Hamed--Bai--Lam canonical forms
\cite{ABL2018}) for all rigorously solved fermionic moment polytopes
(Table~\ref{tab:census}), extending to the pure-state GPC polytopes a bridge
between $N$-representability and positive geometry known previously at the
level of the hypersimplex \cite{Castillo2021,LPW2020}.
\item \textbf{Methods.} An alternating-projection solver for ``pure state
with prescribed occupation spectrum,'' immune to the second-order
degeneracy flatness that stalls gradient and moment-matching methods; a
sparsifying concentration flow, deployed as a jittered multi-start swarm,
which discovered $\psi_A$; and a fully independent verification pipeline,
including a constructive re-verification of all $38$ published extremal
states of $\Pi_{3,8}$.
\end{enumerate}

\section{An explicit extremal state for $v_A$}
\label{sec:theorem}

We write $\ket{i_1 i_2 i_3 i_4}$, $i_1<i_2<i_3<i_4$, for Slater determinants
in $\wedgeN{4}{9}$ over an orthonormal basis of $\mathcal{H}_9$, and call two
determinants \emph{one-hop connected} if they share exactly three orbitals.

\begin{theorem}
\label{thm:psiA}
The normalized state
\begin{equation}
\label{eq:psiA}
\psi_A \;=\; \frac{1}{\sqrt{21}}\Bigl(
\sqrt{2}\,\ket{1257} + \sqrt{3}\,\ket{1347} + \sqrt{2}\,\ket{1369}
+ \sqrt{3}\,\ket{1459} + \sqrt{6}\,\ket{1789} + 2\,\ket{2679} + \ket{3579}
\Bigr)
\end{equation}
has natural occupation numbers exactly
$v_A = \tfrac1{21}(16,16,16,6,6,6,6,6,6)$, with the three heavy natural
orbitals being the basis orbitals $\{1,7,9\}$. Hence the vertex $v_A$ of
$\Pi_{4,9}$ is attained.
\end{theorem}

\begin{proof}
No two of the seven determinants in \eqref{eq:psiA} are one-hop connected
(each pair shares at most two orbitals), so every off-diagonal matrix
element $\langle a_i^\dagger a_j\rangle_{\psi_A}$, $i\neq j$, vanishes
identically: the 1-RDM is diagonal in the given basis, and the amplitudes'
phases are pure gauge. The diagonal entries are the incidence sums
$\lambda_m = \sum_{T \ni m} k_T/21$ with squared weights
$k = (2,3,2,3,6,4,1)$ on the respective determinants. Orbital $1$ lies in
the determinants of weight $2,3,2,3,6$, giving $16/21$; likewise orbitals
$7$ and $9$; each remaining orbital collects weight $6$. Since
$\sum_T k_T = 21$, the state is normalized, and its ordered spectrum is
exactly $v_A$.
\end{proof}

\begin{remark}
Relabeling orbitals so that the heavy modes are $\{1,2,3\}$ gives the
equivalent canonical form
$\psi_A = \tfrac1{\sqrt{21}}(\sqrt6\ket{1239} + \sqrt2\ket{1247}
+ \sqrt3\ket{1256} + \sqrt2\ket{1358} + \sqrt3\ket{1367}
+ 2\ket{2348} + \ket{2357})$.
\end{remark}

\begin{remark}
The proof is checkable by hand in minutes; the discovery was not. The state
was found by the numerical pipeline of Section~\ref{sec:methods} and only
then recognized to have exact structure. We regard the division of labor---
stochastic geometric search for discovery, elementary combinatorics for
proof---as a portable pattern for moment-polytope problems.
\end{remark}

\section{Weighted designs and the dichotomy}

\emph{Terminology and prior art.} The notion below is not new. A support in
which no two determinants differ by a single orbital is, in
representation-theoretic language, a \emph{root-distinct} set of weights: the
weights of the supporting basis vectors differ pairwise by no root of the
acting group. The idea is due to Wildberger~\cite{W92} and was developed by
Sjamaar~\cite{S98}; in the quantum-marginal setting it appears as
Definitions~3--4 of Maciazek and Tsanov~\cite{MT17}, stated there for spinful
and spinless fermions, bosons and distinguishable particles alike, with
applications to qubit systems~\cite{M15} and to pinned fermionic
occupations~\cite{MS20}. Ref.~\cite{MT17} further records that maximal
root-distinct supports were already used to determine the
$\wedge^3\mathcal{H}_6$ and $\wedge^3\mathcal{H}_7$ polytopes---with the
support $\{123,145,246,356\}$ that our census recovers---so our rank-6 and
rank-7 results are a rediscovery, reported here only for completeness. We
retain the term \emph{one-hop independent} for the same condition, since our
arguments are phrased in terms of single-orbital exchanges rather than roots.
What is new here is the algorithmic layer: the exhaustive vertex-by-vertex
decision procedure with solver-independent certificates, the exact repair step
for vertices admitting no such state, and the resulting classification.

\label{sec:census}

\begin{definition}
\label{def:design}
Let $v = n/D$ be a vertex of $\Pi_{N,d}$ with integer numerator vector $n$
and $D = \tfrac1N\sum_m n_m$. A \emph{weighted design} for $v$ is a family
of $N$-subsets $\{T\}$ of $\{1,\dots,d\}$, pairwise sharing at most $N-2$
elements, together with positive real weights $w_T$ with
$\sum_{T\ni m} w_T = n_{\sigma(m)}/D$ for some permutation $\sigma$. The
associated state $\psi = \sum_T \sqrt{w_T}\,\ket{T}$ has exactly diagonal
1-RDM and ordered spectrum $v$; the design is \emph{integer} if
$w_T = k_T/D$ with $k_T \in \mathbb{Z}_{>0}$.
\end{definition}

Design-attainability is thus a feasibility question for a system of integer
incidence equations under an independence (anti-adjacency) condition---no
quantum mechanics remains. We decided it exactly for every vertex of the
rigorously solved small systems, by depth-first search with exact-cover
pivoting and, where search was inconclusive, by mixed-integer linear
programming with binary support variables and pairwise independence
constraints (HiGHS; certificates of infeasibility retained).

\begin{proposition}[Census]
\label{prop:census}
All $4$ vertices of $\Pi_{3,6}$ and all $10$ vertices of $\Pi_{3,7}$ admit
integer weighted designs (reproducing, in the latter case, the extremal
states of Ref.~\cite{AK2008}). Of the $38$ vertices of $\Pi_{3,8}$, exactly
$27$ admit weighted designs ($22$ of them integer at the natural
denominator); the remaining $11$---the exotic vertex
$(5,5,5,5,2,2,2,2)/28$ and its relatives---admit no weighted design with
any positive real weights (mixed-integer certificates with continuous
weight variables and binary support, each independently reproduced with a
second solver of a different lineage, alongside feasibility controls on
design-class vertices). The vertex $v_A$ of $\Pi_{4,9}$
admits the integer design of Theorem~\ref{thm:psiA}.
\end{proposition}

\begin{theorem}[$v_B$ is interference-type]
\label{thm:B}
No weighted design with any positive real weights exists for
$v_B = (20,14,14,14,14,4,4,4,4)/23$. Since the family of Slater determinants
and its one-hop structure are identical in every orthonormal orbital basis,
the statement is basis-free: no extremal state of $v_B$ is one-hop independent in
\emph{any} orbital basis. Equivalently, in any basis in which the 1-RDM of an
extremal state is diagonal, the supporting determinants necessarily fail
one-hop independence; the vertex is attained through interference. (We thank
T.~Maciazek for prompting this basis-free formulation: a $2\times2$
rotation inside the mixing block of the state of Theorem~\ref{thm:psiB}
diagonalizes its 1-RDM while destroying independence, illustrating both halves
of the statement.)
\end{theorem}

\begin{proof}[Proof (computer-assisted)]
A state with one-hop-independent support has exactly diagonal 1-RDM, so
its squared amplitudes would give nonnegative real weights solving the
incidence system on an independent support. The corresponding feasibility
problem---$126$ continuous weight variables, $126$ binary support
variables, $9$ incidence equalities, and $1260$ pairwise independence
inequalities---is infeasible. The certificate was produced by the HiGHS
solver, independently reproduced with the COIN-OR CBC solver, and is
archived with the code; the same certificates were obtained for the
integer-restricted problem at denominators $23$, $46$, and $69$.
\end{proof}


\begin{theorem}[Explicit extremal state for $v_B$]
\label{thm:psiB}
Let $\gamma = \arccos\bigl(3/(4\sqrt{14})\bigr)$. The normalized state
\begin{equation}
\label{eq:psiB}
\psi_B = \tfrac{1}{\sqrt{23}}\Bigl(
2\ket{1236} + \sqrt2\,\ket{1248} + \sqrt7\,e^{i\gamma}\ket{1249}
+ \sqrt3\,\ket{1345} + \sqrt2\,\ket{1378} + \sqrt2\,\ket{1379}
+ \ket{2358} + \sqrt2\,\ket{3489}\Bigr)
\end{equation}
has natural occupation numbers exactly $v_B = \tfrac1{23}(20,14,14,14,14,4,4,4,4)$.
Hence both numerically-verified vertices of Ref.~\cite{AK2008} are attained
exactly, completing the verification of the $\wedge^4\mathcal{H}_9$ moment
polytope.
\end{theorem}

\begin{proof}
The eight determinants pairwise share at most two orbitals except the two
pairs $\{1248,1249\}$ and $\{1378,1379\}$, each differing by the exchange
$8 \leftrightarrow 9$; hence every off-diagonal 1-RDM element vanishes
identically except $\rho_{89}$. The diagonal entries are integer incidence
sums of the squared weights $(4,2,7,3,2,2,1,2)$: orbitals $1$--$7$ receive
$20,14,14,14,4,4,4$ respectively, and orbitals $8,9$ receive $7$ and $11$
(all divided by $23$). The exchange pairs contribute
$\rho_{89} = \bigl(\sqrt{14}\,e^{i\gamma} + 2\bigr)/23
= (11 + i\sqrt{215})/92$ (both fermionic signs are $+1$), so
$|\rho_{89}|^2 = (121+215)/92^2 = 21/23^2$. The remaining $2\times2$ block
$\bigl(\begin{smallmatrix} 7 & z \\ \bar z & 11\end{smallmatrix}\bigr)/23$
has characteristic polynomial $\lambda^2 - 18\lambda + (77 - 21)
= (\lambda-14)(\lambda-4)$, giving eigenvalues $14/23$ and $4/23$ exactly.
The ordered spectrum is $v_B$.
\end{proof}

\begin{theorem}[Real single-block exclusion]
\label{thm:realexclusion}
No state of the form $\psi = \sum_T \varepsilon_T\sqrt{k_T/23}\,\ket{T}$
with $\varepsilon_T \in \{\pm1\}$, positive integers $k_T$, and one-hop
connections confined to a single orbital pair attains $v_B$. (Exhaustive
enumeration: $16$ ansatz equivalence classes, over $4.5\times10^6$
weighted supports, certified optimal terminations.) Consequently the
complex phase in Theorem~\ref{thm:psiB} is not an artifact: within the
single-block class, $v_B$ is attainable only with nontrivial relative phase.
\end{theorem}

\begin{remark}[Non-realifiability under orbital rotations (numerical)]
\label{rem:orbit}
Coefficient complexity is basis-dependent, so a sharper question is whether
any orbital rotation makes the state of Theorem~\ref{thm:psiB} real. A state
$\psi$ admits a real form iff there exists $V \in U(9)$ with
$\Lambda^4 V\,\bar\psi = e^{i\alpha}\psi$; such $V$ must map the
eigenspaces of $\bar\rho$ to those of $\rho$, reducing the search to
$U(1)\times U(4)\times U(4)$ aligned with the natural eigenspaces. Diagonal
(phase-only) rotations are excluded analytically: the two exchange pairs force
contradictory phase differences $2\gamma$ and $0$ on the same mode pair.
Numerically, maximizing $|\langle\psi,\Lambda^4 V\bar\psi\rangle|$ over
the full reduced group converges from independent random starts to the same
value $\approx 0.933184 < 1$, indicating that $\psi_B$ admits no real form
under any orbital rotation, with the deficit an apparent invariant of its
orbit. This is a statement about the constructed state, not the vertex:
whether $v_B$ admits some other, real extremal state remains open (exhaustive
enumeration excludes integer-weight real single-exchange-block states; the
continuous-weight real case on larger supports, and deeper block structures,
remain open).
\end{remark}

\begin{remark}[Relation to the superselection rule of Ref.~\cite{Liebert2025}]
Liebert \emph{et al.}\ prove that saturation of a (spin-adapted) GPC
restricts which configurations may contribute to the wave function.
A vertex saturates many GPCs at once; the dichotomy sharpens the picture at
this extreme point: for design vertices the intersected selection rules
admit a diagonal, effectively classical weighted configuration mixture,
while for interference vertices they force coherent superpositions of
connected configurations. The historically difficult vertices of
Ref.~\cite{AK2008} are exactly the interference ones, which we propose as
the structural reason they evaded sparse representation-theoretic
constructions.
\end{remark}

\section{GPC polytopes as positive geometries}
\label{sec:posgeo}

Every convex polytope carries a canonical form in the sense of
Arkani-Hamed--Bai--Lam \cite{ABL2018}: a rational top-form with simple poles
exactly on facets, whose numerator is the adjoint polynomial
\cite{KohnRanestad}. A connection between $N$-representability and positive
geometry was observed at the level of the hypersimplex
\cite{Castillo2021,LPW2020}; the pure-state GPC polytopes themselves appear
not to have been examined in this light. For the rigorously solved systems the invariants are: $\Pi_{3,6}$ ($4$ vertices, $4$ facets, adjoint degree $0$; a simplex), $\Pi_{3,7}$ ($10$, $10$, $3$), $\Pi_{3,8}$ ($38$, $39$, $31$), and $\Pi_{4,8}$ ($22$, $23$, $15$), with adjoint degree $=F-\dim-1$. Table~\ref{tab:census} reports the
basic invariants for all rigorously solved systems, computed from exact
rational vertex/facet enumeration (dual-volume evaluations of the canonical
forms and numerical residue tests on facets---including the exotic facets of
$\Pi_{3,8}$---are described in the repository).


Two structural observations from the same computations: positivity of the
smallest occupation number is \emph{implied} by the GPCs at ranks $6$ and
$7$ and returns as a genuine facet at rank $8$; and both rank-$8$ systems
satisfy $F = V+1$, which we flag without explanation.

\section{Methods}
\label{sec:methods}

\paragraph{Alternating projections for prescribed spectra.}
Optimizing any functional of the eigenvalues of the 1-RDM stalls near
degenerate targets: eigenvalues respond only at second order to
intra-cluster perturbations, so matching moments to machine precision
$\varepsilon$ constrains the spectrum only to
$O(\sqrt{\varepsilon})$---which quantitatively reproduces the
$\sim\!10^{-8}$ floors we observed for gradient, moment-matching (L-BFGS
with verified analytic gradients), and support-function methods. The cure is
to keep eigenvalues out of the inner loop: alternately (i) fix the natural
orbitals $U$ of the current state and set the target matrix
$T = U\,\mathrm{diag}(v)\,U^\dagger$, then (ii) minimize the smooth
quartic $\|\rho(\psi) - T\|_F^2$ by L-BFGS with exact gradients. The
resulting solver reaches prescribed spectra to $10^{-9}$ from random
initializations, and its calibration on vertices with known extremal states
gated every claim in this paper.

\paragraph{Concentration flow and swarm discovery.}
On the solution manifold $\{\psi : \operatorname{spec}\rho(\psi) = v\}$
(dimension $\approx 242$ for $v_A$) we ran a sparsifying flow: ascent of the
inverse participation ratio $\sum_T |c_T|^4$, alternated with full-power
reprojection, guarded by rejecting any step whose reprojected spectral
distance exceeds $10^{-8}$, with hard removal of sub-threshold amplitudes
under the same guard. Deployed as a multi-start swarm (12 jittered walkers,
consumer hardware), one walker cascaded from $126$ to $7$ determinants at
spectral distance $2.3\times10^{-14}$ within minutes; the resulting support
and rational squared amplitudes were then recognized and proved exactly
(Theorem~\ref{thm:psiA}). We note for honesty that single-trajectory runs of
the same flow plateau near support $120$: conclusions about
$242$-dimensional solution manifolds require swarm-scale exploration, and an
early version of this project wrongly concluded from local searches that the
extremal states were irreducibly dense.

\paragraph{Verification pipeline.}
All claims passed: (i) independent re-derivation of the 1-RDM with separate
fermionic-sign bookkeeping; (ii) exact rational arithmetic for the final
proofs (no floating point); (iii) constructive re-verification of all $38$
published extremal states of $\Pi_{3,8}$ from Table~6 of
Ref.~\cite{AK2008}, each reproducing its claimed vertex to machine
precision; and (iv) calibration of every attainability method on vertices
with known extremal states before its application to open ones.

\section{Scaling toward chemically relevant dimensions}
\label{sec:scaling}
The gap this paper addresses at rank $9$ is an instance of a larger one:
constraint tables exist only for $d \le 10$ (with $(3,d)$ systems reported to
$d = 12$), while quantum-chemical applications need $d \sim 20$--$40$; the
polytope program has largely paused at this frontier, in part because
extremal-state construction has been manual. We outline a four-instrument
program toward that regime, distinguishing carefully between rigorous and
heuristic layers.

\emph{(i) Exact inner structure (rigorous; this work, scaled).} Root-distinct
enumeration produces certified rational polytope points at any rank: the
CP-SAT/MILP feasibility problems grow with $\binom{d}{N}$---$1140$
determinants at $(3,20)$, $9880$ at $(3,40)$---within modern solver reach,
and every constructed state is an exact extreme-point candidate with a
hand-checkable proof. The exchange-block repair step is a per-candidate
polygon test whose cost is independent of $d$. This yields exact inner
approximations (design hulls and their repairs) at all chemically relevant
dimensions on commodity hardware.

\emph{(ii) Rigorous boundary anchors from known ranks.} The polytope
$\Delta(N,d)$ is precisely the face $\{\lambda_{d+1}=0\}$ of
$\Delta(N,d+1)$, so the rigorously known tables at $d \le 12$ supply exact
face data---boundary anchors---for every higher rank, and the closed-form
Grassmann inequality families remain valid at all $d$. (We note that valid
inequalities do \emph{not} in general lift from $\Delta(N,d)$ to
$\Delta(N,d+1)$; the face relation, not naive lifting, is the correct
rigorous statement.)

\emph{(iii) Pointwise membership oracles (rigorous, pointwise).} Membership
in moment polytopes of reductive group representations---including the
fermionic case explicitly---is in $\mathrm{NP}\cap\mathrm{coNP}$
\cite{BCMW17}, and the tensor-scaling algorithms of
Ref.~\cite{BFGOWW18} provide an efficient weak membership oracle. Since the
oracle iterates on state vectors of dimension $\binom{d}{N}$, pointwise
boundary location (bisection along rays) is computationally cheap even at
$(3,40)$. Our alternating-projection solver is a primitive relative of these
methods; the scaling literature supplies convergence guarantees we lacked.

\emph{(iv) Toward full outer certification.} The representation-theoretic
computation of complete constraint lists becomes prohibitive beyond
$d \approx 12$ along the original route. A recent algorithmic line
\cite{vdBerg25} computes moment polytopes for general reductive-group
representations via Franz's description, with a current frontier of
$27$--$64$-dimensional representations (all of
$\mathbb{C}^3{\otimes}\mathbb{C}^3{\otimes}\mathbb{C}^3$, and
$\mathbb{C}^4{\otimes}\mathbb{C}^4{\otimes}\mathbb{C}^4$ with high
probability). This does not yet reach the fermionic tables (already at
representation dimensions $126$--$252$ in 2008), but it scales along a
different complexity axis than the classical approach, and to our knowledge
has not been applied to $\Lambda^N\mathbb{C}^d$. Validating it against the
known fermionic tables, and attempting $(3,13)$---which would be the first
new generalized Pauli constraint system since 2008---is a concrete
experiment we highlight.

\emph{Resources.} Instruments (i)--(iii) run at $(3,d)$ for $d \le 20$ on a
single workstation in weeks; extending the inner atlas to $(3,40)$ is a
modest cloud campaign (order $10^2$ vCPU-months). Instrument (iv) is
gated on algorithm engineering rather than hardware. The binding constraint
throughout is machine-readable vertex/facet data for the known systems
(ranks $9$--$12$), for validation; we would welcome pointers to surviving
copies of the original datasets.

\section{The census at ranks 9 and 10}
\label{sec:census910}

\begin{table}[t]
\centering
\begin{tabular}{lcccccc}
\toprule
polytope & $\dim$ & vertices & constraints & design (int) & design (real) & interference \\
\midrule
$\Pi_{3,6}$ & 3 & 4 & $1{+}3$eq & \multicolumn{3}{c}{all design (doubly-excited regime \cite{MT17})} \\
$\Pi_{3,7}$ & 6 & 10 & 4 & \multicolumn{3}{c}{all design (doubly-excited regime \cite{MT17})} \\
$\Pi_{3,8}$ & 7 & 38 & 31 & 26 & 1 & 11 \\
$\Pi_{4,8}$ & 7 & 22 & 15 & 22 & 0 & 0 \\
$\Pi_{3,9}$ & 8 & 58 & 52 & 37 & 1 & 20 \\
$\Pi_{4,9}$ & 8 & 103 & 60 & 85 & 2 & 16 \\
$\Pi_{3,10}$ & 9 & 113 & 93 & 69 & 2 & 42 \\
$\Pi_{4,10}$ & 9 & 159 & 125 & 132 & 2 & 25 \\
$\Pi_{5,10}$ & 9 & 292 & 161 & 246 & 4 & 42 \\
\bottomrule
\end{tabular}
\caption{Every determinate fermionic moment polytope, with the
design/interference census. Ranks $\le 8$ are rigorously complete; the
rank-$10$ constraint lists are conjecturally complete as in
Ref.~\cite{AK2008}, and all rank-$10$ rows are conditional on them.
Verdicts transport under particle--hole duality, extending the
classification to all $(N,d)$ with $d\le 10$. For $\Pi_{4,8}$ (half
filling) the particle--hole involution maps vertex and facet sets to
themselves; by uniqueness of the canonical form, $\Omega$ is exactly
PH-invariant. Interference is absent at rank $8$ in the $N{=}4$ series
and first appears at rank $9$---the system of this paper. The $16$
interference vertices of $\wedge^4\mathcal{H}_9$ (integer forms at
natural denominator):
$(12{,}9{,}5{,}5{,}5{,}3{,}3{,}3{,}3)$,
$(10{,}7{,}7{,}4{,}4{,}4{,}2{,}1{,}1)$,
$(10{,}7{,}5{,}5{,}5{,}2{,}2{,}2{,}2)$,
$(9{,}6{,}6{,}4{,}4{,}4{,}1{,}1{,}1)$,
$(18{,}12{,}12{,}7{,}7{,}4{,}4{,}4{,}4)$,
$(9{,}6{,}5{,}5{,}5{,}2{,}2{,}1{,}1)$,
$(15{,}9{,}8{,}8{,}8{,}3{,}3{,}3{,}3)$,
$(7{,}4{,}4{,}4{,}4{,}2{,}1{,}1{,}1)$,
$(16{,}9{,}9{,}9{,}9{,}4{,}4{,}2{,}2)$,
$(18{,}10{,}10{,}10{,}10{,}4{,}4{,}3{,}3)$,
$(28{,}15{,}15{,}15{,}15{,}6{,}6{,}6{,}6)$,
$(20{,}12{,}12{,}12{,}12{,}4{,}4{,}4{,}4)$,
$(14{,}11{,}11{,}6{,}6{,}6{,}2{,}2{,}2)$,
$(14{,}9{,}9{,}9{,}9{,}3{,}3{,}2{,}2)$,
$(11{,}8{,}7{,}7{,}7{,}2{,}2{,}2{,}2)$,
and $v_B=(20{,}14{,}14{,}14{,}14{,}4{,}4{,}4{,}4)$.
Real-not-integer vertices: $(6{,}3{,}3{,}3{,}3{,}3{,}1{,}1{,}1)$ and
$(10{,}10{,}10{,}7{,}7{,}2{,}2{,}2{,}2)$. Full lists, verdicts, and
validation code accompany the data bundle.}
\label{tab:census}
\end{table}



After the results above were obtained, we carried out the program proposed in
the Discussion: extending the design/interference census to every rank-$9$ and
rank-$10$ system. The machine-readable constraint data of Ref.~\cite{AK2008}
(hosted at Bilkent) is no longer accessible; we recovered the complete
inequality lists for $\wedge^3\mathcal{H}_9$, $\wedge^4\mathcal{H}_9$,
$\wedge^3\mathcal{H}_{10}$, $\wedge^4\mathcal{H}_{10}$ and
$\wedge^5\mathcal{H}_{10}$ from the text of Altunbulak's thesis
\cite{AltunbulakThesis}, with parsed counts matching the published totals
($52$, $60$, $93$, $125$, $161$). Exact vertex enumeration (lrs, rational
arithmetic) and the MILP census of Section~5 were then applied to all five
systems, with every step cross-validated by independent structural
invariants: embedding coherence ($\Delta(N,d)$ is the face
$\{\lambda_{d+1}=0\}$ of $\Delta(N,d+1)$), frozen-core lifts
$\lambda \mapsto (1,\lambda)$, particle--hole self-duality of
$\wedge^5\mathcal{H}_{10}$ at the level of inequalities, vertices, and
verdicts, and agreement with the published vertex tables at rank $8$.

(One inequality of the $\wedge^3\mathcal{H}_9$ table of Ref.~\cite{AltunbulakThesis} is lost at a page boundary; we work with the repaired $52$-inequality system, which yields $58$ vertices consistent with all cross-checks. Details accompany the data bundle.)

\textbf{The census.} Table~\ref{tab:census} summarizes the classification.
For the system of this paper, $\wedge^4\mathcal{H}_9$: of $103$ vertices
(matching the count of Ref.~\cite{AK2008}), $85$ admit integer weighted
designs at the natural denominator, $2$ admit real but not integer designs,
and $16$ require interference. The vertex $v_A$ is classified DESIGN
(consistent with Theorem~1) and $v_B$ INTERFERENCE (consistent with
Theorem~2), each verdict landing blind inside a uniform sweep. Notably,
$v_B$ is not isolated: the vertex
$(20{:}12{:}12{:}12{:}12{:}4{:}4{:}4{:}4)/21$ shares its architecture
(head $20$, degenerate quadruple, tail $4^4$) at the neighbouring natural
denominator and likewise requires interference, so $v_B$ is the first-found
member of a family. Across ranks $9\to 10$ the interference fraction is stable within each particle-number series ($34\to 37\%$ at $N{=}3$, $16\%$ at $N{=}4$, $14\%$ at $N{=}5$; the rank-$8$ systems sit lower, $29\%$ and $0\%$) while padding and frozen-core
lifting generate the majority of each rank's interference vertices from
lower-rank originals; at $\wedge^4\mathcal{H}_{10}$ they generate
\emph{all} of them. Verdicts transport under particle--hole duality (the
complement bijection on determinants maps designs to designs at the same
denominator), so the classification extends to all systems
$(N,d)$, $d \le 10$. The rank-$10$ constraint lists remain conjecturally
complete as in Ref.~\cite{AK2008}; all rank-$10$ statements here are
conditional on those lists, which our tests probe but cannot fully certify.


\section{Discussion}

The remaining open object is an explicit extremal state for $v_B$, now known
to require interference; a natural attack is exact phase-cancellation
solving over small connected supports, guided by the MILP census, or the
facet conditions of Ressayre-type inequalities \cite{Ressayre}. The
dichotomy itself invites theory: which incidence data of moment-polytope
vertices are weighted-independent-set realizable is a clean combinatorial
question, and the census suggests the answer encodes the boundary of
``classically representable'' vertices. The entire toolkit transfers to
other pure-state quantum marginal settings---most directly the entanglement
polytopes of Ref.~\cite{WDGC2013}---where the same questions (explicit
vertex states, design versus interference, canonical forms) appear open.
Section~\ref{sec:census910} carries out this extension: the census now covers every rank-$9$ and rank-$10$ system, cross-validated by embedding, lifting, and duality invariants.


\subsection*{Three questions this raises}
(i) \emph{Does the complex class exist, and persist?} Theorem~\ref{thm:psiB}
exhibits the first genuinely complex extremal state known for these polytopes,
Theorem~\ref{thm:realexclusion} shows real amplitudes cannot be repaired
within its mechanism class, and Remark~\ref{rem:orbit} indicates the state
admits no real form in any orbital basis; whether $v_B$ admits \emph{any}
real extremal state is open, as is the same question for the rank-$9$ and rank-$10$ polytopes, whose interference vertices Section~\ref{sec:census910} now enumerates ($16$ at $\wedge^4\mathcal{H}_9$ alone, including a structural sibling of $v_B$); whether each admits a real extremal state, and whether field of definition is decidable vertex-by-vertex the way design-attainability is, remain open. We note that numerical verification cannot
distinguish the classes: real-restricted optimization approaches $v_B$ to
working precision even though no real state attains it. (ii) \emph{Is the
design/real/complex trichotomy an instance of a broader hierarchy?} The three
classes are naturally reminiscent of the positive/signed/complex gradation
familiar from classical simulability and sign-problem considerations; whether
that resemblance is structural or superficial we leave open. (iii)
\emph{What is the physics of a pinned phase?} Occupation spectra are
phase-blind, yet pinning to $v_B$ forces the relative phase
$\cos\gamma = 3/(4\sqrt{14})$ exactly; the dynamical and metrological
consequences of a geometrically pinned algebraic phase appear unexplored.

\section*{Acknowledgments}
We are grateful to Tomasz Maciazek for generous and rapid
correspondence on a first version of this manuscript: for pointing out the
root-distinct literature and lineage now cited in
the weighted-designs section, for the basis-free reformulation prompt
absorbed into Theorem~\ref{thm:B}, and for emphasizing that the algorithmic
construction, rather than any particular state, is the contribution most
likely to be useful---advice that reshaped this draft.
The computations, literature analysis, and drafting in this work were
carried out in close collaboration with Claude (Anthropic); all results
were verified through the independent pipeline described in
Section~\ref{sec:methods}, and responsibility for the claims rests with the
author. The multi-start swarm computations were performed on the author's
personal hardware.

\begin{thebibliography}{99}
\bibitem{AltunbulakThesis} M.~Altunbulak, \emph{The Pauli principle, representation theory, and geometry of flag varieties}, Ph.D.\ thesis, Bilkent University (2008).
\bibitem{Klyachko2006} A.~Klyachko, \emph{Quantum marginal problem and
$N$-representability}, J.\ Phys.: Conf.\ Ser.\ \textbf{36}, 72 (2006).
\bibitem{BCMW17} P.~B\"urgisser, M.~Christandl, K.~D.~Mulmuley, and
M.~Walter, \emph{Membership in moment polytopes is in NP and coNP},
SIAM J.~Comput.\ \textbf{46}, 972 (2017), arXiv:1511.03675.
\bibitem{BFGOWW18} P.~B\"urgisser, C.~Franks, A.~Garg, R.~Oliveira,
M.~Walter, and A.~Wigderson, \emph{Efficient algorithms for tensor scaling,
quantum marginals, and moment polytopes}, FOCS 2018, arXiv:1804.04739.
\bibitem{vdBerg25} M.~van den Berg, M.~Christandl, V.~Lysikov,
H.~Nieuwboer, M.~Walter, and J.~Zuiddam, \emph{Computing moment
polytopes---with a focus on tensors, entanglement and matrix
multiplication}, arXiv:2510.08336; STOC 2025.
\bibitem{W92} N.~J.~Wildberger, \emph{The moment map of a Lie group
representation}, Trans.\ Amer.\ Math.\ Soc.\ \textbf{330}, 257 (1992).
\bibitem{S98} R.~Sjamaar, \emph{Convexity properties of the moment mapping
re-examined}, Adv.\ Math.\ \textbf{138}, 46 (1998).
\bibitem{MT17} T.~Maciazek and V.~Tsanov,
\emph{Quantum marginals from pure doubly excited states},
J.~Phys.~A: Math.~Theor.\ \textbf{50}, 465304 (2017),
doi:10.1088/1751-8121/aa8c5f.
\bibitem{M15} T.~Maciazek and A.~Sawicki,
\emph{Critical points of the linear entropy for pure $L$-qubit states},
J.~Phys.~A: Math.~Theor.\ \textbf{48}, 045305 (2015),
doi:10.1088/1751-8113/48/4/045305.
\bibitem{MS20} T.~Maciazek, A.~Sawicki, D.~Gross, A.~Lopes, and
C.~Schilling, \emph{Implications of pinned occupation numbers for natural
orbital expansions. II: rigorous derivation and extension to non-fermionic
systems}, New J.~Phys.\ (2020), doi:10.1088/1367-2630/ab64b1.
\bibitem{AK2008} M.~Altunbulak and A.~Klyachko, \emph{The Pauli principle
revisited}, Comm.\ Math.\ Phys.\ \textbf{282}, 287--322 (2008).
\bibitem{BD1972} R.~E.~Borland and K.~Dennis, \emph{The conditions on the
one-matrix for three-body fermion wavefunctions with one-rank equal to six},
J.\ Phys.\ B \textbf{5}, 7 (1972).
\bibitem{SGC2013} C.~Schilling, D.~Gross, and M.~Christandl, \emph{Pinning
of fermionic occupation numbers}, Phys.\ Rev.\ Lett.\ \textbf{110}, 040404
(2013).
\bibitem{Liebert2025} J.~Liebert, Y.~Lemke, M.~Altunbulak, T.~Maciazek,
C.~Ochsenfeld, and C.~Schilling, \emph{Toolbox of spin-adapted generalized
Pauli constraints}, Phys.\ Rev.\ Research \textbf{7}, 023247 (2025).
\bibitem{Castillo2021} F.~Castillo, J.-P.~Labb\'e, J.~Liebert, A.~Padrol,
E.~Philippe, and C.~Schilling, \emph{An effective solution to convex
1-body $N$-representability}, arXiv:2105.06459.
\bibitem{ABL2018} N.~Arkani-Hamed, Y.~Bai, and T.~Lam, \emph{Positive
geometries and canonical forms}, JHEP \textbf{11} (2017) 039.
\bibitem{KohnRanestad} K.~Kohn and K.~Ranestad, \emph{Projective geometry
of Wachspress coordinates}, Found.\ Comput.\ Math.\ \textbf{20}, 1135
(2020).
\bibitem{LPW2020} T.~Lukowski, M.~Parisi, and L.~K.~Williams, \emph{The
positive tropical Grassmannian, the hypersimplex, and the m=2
amplituhedron}, arXiv:2002.06164.
\bibitem{Ressayre} N.~Ressayre, \emph{Geometric invariant theory and the
generalized eigenvalue problem}, Invent.\ Math.\ \textbf{180}, 389 (2010).
\bibitem{WDGC2013} M.~Walter, B.~Doran, D.~Gross, and M.~Christandl,
\emph{Entanglement polytopes: multiparticle entanglement from
single-particle information}, Science \textbf{340}, 1205 (2013).
\end{thebibliography}

\end{document}
